\documentclass{article}
\usepackage[margin=1in]{geometry}
\usepackage{amsmath,amssymb,amsthm}
\usepackage{array,booktabs,longtable}
\usepackage{hyperref}

\newtheorem{definition}{Definition}
\newtheorem{theorem}[definition]{Theorem}
\newtheorem{lemma}[definition]{Lemma}
\newtheorem{proposition}[definition]{Proposition}
\newtheorem{corollary}[definition]{Corollary}
\newtheorem{convention}[definition]{Convention}

\newcommand{\source}[1]{\par\noindent\textbf{Source.} #1}
\newcommand{\lean}[1]{\par\noindent\textbf{Lean target.} \texttt{\detokenize{#1}}}
\newcommand{\status}[1]{\par\noindent\textbf{Status.} #1}
\newcommand{\depends}[1]{\par\noindent\textbf{Depends on.} #1}
\newcommand{\backend}[1]{\par\noindent\textbf{Lean backend.} #1}
\newcommand{\cA}{\mathcal A}
\newcommand{\cH}{\mathcal H}
\newcommand{\cM}{\mathcal M}
\newcommand{\Ftwo}{\mathbb F_2}

\begin{document}

\section*{Blueprint 006: Full Type-Classification Dependency Plan}

\subsection*{Scope}

This master blueprint supports the strengthened goal: prepare all theory
targets needed before Lean formalization of the full canonical stabilizer
projection-state von Neumann algebra type classification.  It is not a Lean
file and it is not an Aristotle packet.  It is the manifest-level dependency
plan that fixes the target IDs, the mathematical order, the expected Lean
backend, and the source locators for the operator-algebra results consumed by
the final theorem.

\paragraph{Historical status.}
This file is now a preserved dependency-plan ledger, not an active Aristotle
source.  The active per-target blueprint notes are
\path{blueprints/007-pauli-quasilocal-algebra.tex} through
\path{blueprints/013-final-stabilizer-type-classification.tex}, whose target
environments use the current style
\(\texttt{LeanFacingName, ID: TARGET, source locator}\).  Do not extract new
packets directly from this ledger until any stale target environment here has
been normalized against the active notes.

The final theorem to be formalized is the classification of the represented
region von Neumann algebra
\[
  M_\Gamma(\omega)=\pi_\omega(\cA_\Gamma)''
\]
for canonical stabilizer projection states on finite or countably infinite
regions.  The classification is by the internal fixed rank \(p\) and internal
defect rank \(r\):
\[
\begin{array}{c|c|c}
p<\infty & r<\infty & \hbox{finite type } I\\
p=\infty & r<\infty & \hbox{type } I_\infty\\
p<\infty & r=\infty & \hbox{hyperfinite type } II_1\\
p=\infty & r=\infty & \hbox{type } II_\infty.
\end{array}
\]
No type \(III\) algebra occurs in this canonical stabilizer sector.

\subsection*{Conventions}

\begin{convention}[Ambient conventions]
All Hilbert spaces are complex Hilbert spaces.  A represented von Neumann
algebra is a concrete star subalgebra of bounded operators on a Hilbert space,
equal to its bicommutant.  Infinite regions are finite or countably infinite
sets of qubit modes.  The Pauli label space of a region \(\Gamma\) is the
finite-support binary symplectic vector space
\[
  V_\Gamma=(\Gamma\to_0\Ftwo)\times(\Gamma\to_0\Ftwo).
\]
The algebraic symplectic form is the standard Pauli commutation form.
\source{Gottesman, \emph{Surviving as a Quantum Computer in a Classical World},
Chapter 3.5, Definitions 3.11 and 3.13, PDF pp. 55--56; project proof note
\path{runs/2026-06-18-stabilizer-algebra-classification/artifacts/004-internal-stabilizer-reduction-type-classification-R1.tex},
Sections 2--5.}
\status{Convention for all targets in this file.}
\end{convention}

\begin{convention}[Mathlib reliance]
The Lean development should rely on mathlib where the coverage audit found
usable declarations: Hilbert spaces, bounded operators, positive linear maps,
GNS via \path{PositiveLinearMap.GNS}, concrete
\path{VonNeumannAlgebra H}, commutants, star projections, finite matrix
traces, \(\texttt{CStarMatrix}\), and finite Hilbert tensor products.  The
project should not assume mathlib already contains factor/type classification,
Murray--von Neumann comparison, semifinite trace theory, or infinite tensor
product von Neumann algebra theory.
\source{\path{agents/2026-06-21-mathlib-operator-algebra-coverage-audit.md}.}
\status{Mathlib-backed foundation rule.}
\end{convention}

\subsection*{Source Map}

\begin{longtable}{@{}p{0.26\linewidth}p{0.66\linewidth}@{}}
\toprule
Source & Locator used in this blueprint \\
\midrule
\endhead
Gottesman QECC notes & Chapter 3.5, Definitions 3.11 and 3.13, PDF
pp. 55--56, for stabilizer groups, Pauli labels, and stabilizer-state
expectation conventions. \\
Project proof note 004 & Sections 2--5 and Theorem 5.1 for internal
stabilizers, crossing-stabilizer invisibility, adapted symplectic bases,
product form, and the stabilizer-side type table. \\
Kadison--Ringrose I & Theorem 4.5.2 and Proposition 4.5.3, PDF pp. 287--288,
for the GNS construction and uniqueness. \\
Kadison--Ringrose II & Definition 6.5.1, Theorem 6.5.2, Corollary 6.5.3,
Remark 6.5.4, and Theorem 6.6.1, PDF pp. 57--63, for type decomposition,
factor type alternatives, type \(I_n\) factors as \(B(\cH)\), and the type
III criterion. \\
Blackadar & III.3.1.1--III.3.1.6, PDF pp. 301--304, for incomplete infinite
tensor products, product-state GNS representations, factor preservation,
the tracial \(M_2\) product \(R\), and \(R\overline\otimes B(\cH)\) as a
type \(II_\infty\) factor. \\
Blackadar & III.3.4.1--III.3.4.2, PDF p. 317, for hyperfinite von Neumann
algebras and infinite tensor products of finite-dimensional algebras being
approximately finite dimensional. \\
\bottomrule
\end{longtable}

\subsection*{Target Dependency Table}

\begin{longtable}{@{}p{0.13\linewidth}p{0.29\linewidth}p{0.18\linewidth}p{0.30\linewidth}@{}}
\toprule
ID & Target & Backend & Main dependencies \\
\midrule
\endhead
SYS-001 & Algebraic stabilizer system & done / project-local & source ledger \\
PAULI-001 & Pauli labels and commutation & project-local, mathlib finite support & SYS-001 \\
PAULI-002 & Finite Pauli algebra and trace basis & project-local, mathlib matrices & PAULI-001 \\
QL-001 & Quasi-local qubit algebra/local net & project-local abstraction & PAULI-002 \\
QL-002 & Region algebras and support restriction & project-local abstraction & QL-001 \\
STATE-001 & Canonical coefficient rule & done / project-local & SYS-001 \\
STATE-002 & Finite-local positivity and consistency & project-local & PAULI-002, STATE-001 \\
STATE-003 & Quasi-local state extension & project-local, mathlib positive maps & STATE-002, QL-001 \\
STATE-004 & Region restriction and crossing invisibility & project-local & QL-002, STATE-003 \\
SYMP-001 & Adapted countable symplectic basis & project-local algebra & PAULI-001 \\
SYMP-002 & Adapted basis gives Pauli algebra isomorphism & project-local & SYMP-001, QL-001 \\
PROD-001 & Product form of the state & project-local & STATE-004, SYMP-002 \\
HILB-001 & Infinite Hilbert-space interface & mathlib-backed plus project wrappers & STATE-003 \\
GNS-001 & GNS package for states & mathlib-backed wrapper & HILB-001, STATE-003 \\
GNS-002 & GNS unitary equivalence under state transport & project-local wrapper & GNS-001, SYMP-002 \\
GNS-003 & Product-state GNS identification & project-local theorem & PROD-001, GNS-002 \\
VN-001 & Represented region von Neumann algebra & mathlib-backed definition & GNS-001, QL-002 \\
VN-002 & Centers and factor predicate & project-local & VN-001 \\
VN-003 & Bicommutants under unitary equivalence & mathlib-backed plus wrappers & VN-001, GNS-002 \\
PROJ-001 & Projection equivalence and finiteness & project-local VN theory & VN-002 \\
TRACE-001 & Finite and semifinite traces & project-local VN theory & PROJ-001 \\
TYPE-DEF-001 & Type \(I_n\), \(I_\infty\), \(II_1\), \(II_\infty\), \(III\) & project-local VN theory & PROJ-001, TRACE-001 \\
MATRIX-001 & Matrix and \(B(\ell^2)\) model factors & project-local, mathlib matrices & TYPE-DEF-001 \\
ITP-001 & Infinite tensor product VN algebra package & project-local sourced theorem & GNS-003, TYPE-DEF-001 \\
ITP-PURE-001 & Infinite pure-mode product is type \(I\) & project-local sourced theorem & ITP-001, MATRIX-001 \\
ITP-TRACE-001 & Infinite tracial defect product is hyperfinite \(II_1\) & project-local sourced theorem & ITP-001, MATRIX-001 \\
TENSOR-001 & Tensor-product type calculus for the four factors & project-local sourced theorem & ITP-PURE-001, ITP-TRACE-001 \\
TYPE-001 & Full stabilizer projection-state type classification & final project theorem & PROD-001, VN-003, TENSOR-001 \\
TYPE-002 & Type \(III\) exclusion corollary & final project corollary & TYPE-001, TYPE-DEF-001 \\
\bottomrule
\end{longtable}

\section{Stabilizer and Quasi-Local Algebra Layer}

\begin{definition}[ID: PAULI-001]
For a finite or countably infinite region \(\Gamma\), the Pauli label space is
\[
  V_\Gamma=(\Gamma\to_0\Ftwo)\times(\Gamma\to_0\Ftwo).
\]
For \(v=(x,z)\) and \(w=(x',z')\), define the binary symplectic commutation
form
\[
  \langle v,w\rangle
  =
  \sum_{j\in\Gamma} x_j z'_j+\sum_{j\in\Gamma} z_j x'_j
  \in \Ftwo,
\]
where both sums are finite because labels have finite support.  Two Pauli
representatives commute exactly when this value is zero.
\source{Gottesman, Chapter 3.5, Definitions 3.11 and 3.13, PDF pp. 55--56;
project proof note 004, Section 2.}
\lean{Formalization.Stabilizer.PauliLabel; Formalization.Stabilizer.symplecticForm}
\backend{Mostly already present through SYS-001; later targets need the
commutation theorem connecting labels to finite matrix representatives.}
\status{Project-local blueprint target, partly implemented by SYS-001.}
\end{definition}

\begin{lemma}[ID: PAULI-002]
For every finite \(F\subset \Gamma\), the Pauli representatives indexed by
\(V_F\) form a linear basis of the matrix algebra
\[
  \cA_F\cong \bigotimes_{j\in F} M_2(\mathbb C).
\]
With the normalized finite trace \(\operatorname{tr}_F\), the Pauli basis is
orthogonal and the coefficient of a local observable \(a\in\cA_F\) at a label
\(v\in V_F\) is \(\operatorname{tr}_F(P_v^*a)\).
\source{Finite-dimensional Pauli calculus internal to project proof note 004,
Sections 2 and 3; Gottesman, Chapter 3.5, PDF pp. 55--56, for Pauli label
conventions.}
\lean{Formalization.Stabilizer.FinitePauliAlgebra}
\backend{Use mathlib finite matrices, \(\texttt{CStarMatrix}\), matrix trace,
and tensor/Kronecker trace facts where possible.}
\status{Planned project-local blueprint target.}
\end{lemma}

\begin{proof}[Proof sketch]
For a single qubit, \(I,X,Z,XZ\) are four linearly independent \(2\times2\)
matrices and hence a basis of \(M_2(\mathbb C)\).  Tensor products over a
finite set \(F\) give \(4^{|F|}\) linearly independent matrices, matching the
dimension of \(\bigotimes_{j\in F}M_2(\mathbb C)\).  Orthogonality follows
from the one-site normalized trace identities and multiplicativity of the
finite tensor trace.
\end{proof}

\begin{definition}[ID: QL-001]
The qubit quasi-local algebra over \(\Gamma\) is the norm closure of the
directed union of finite local algebras:
\[
  \cA_\Gamma=\overline{\bigcup_{F\Subset\Gamma}\cA_F}.
\]
For a Lean-first development, this may initially be represented as a local-net
abstraction consisting of finite algebras \(\cA_F\), inclusion maps
\(\cA_F\to\cA_G\) for \(F\subset G\), and a universal state/restriction
interface sufficient for GNS and region algebras.
\source{Project proof note 004, Section 2; Blackadar, general C*-tensor
product background around III.3.1.2, PDF pp. 301--302, for infinite products
as represented closures.}
\lean{Formalization.Stabilizer.QuasiLocalAlgebra}
\backend{Project-local abstraction; use mathlib C*-algebra classes after the
algebra carrier is chosen.}
\status{Planned project-local blueprint target.}
\end{definition}

\begin{definition}[ID: QL-002]
For a region \(\Lambda\subseteq\Gamma\), the region algebra
\(\cA_\Lambda\subseteq\cA_\Gamma\) is the C*-subalgebra generated by finite
local algebras supported in finite \(F\subset\Lambda\).  Its finite-support
Pauli labels are exactly \(V_\Lambda\subseteq V_\Gamma\).
\source{Project proof note 004, Section 2.}
\lean{Formalization.Stabilizer.RegionAlgebra}
\backend{Project-local abstraction, later connected to mathlib subalgebras.}
\status{Planned project-local blueprint target.}
\end{definition}

\section{Canonical States}

\begin{theorem}[ID: STATE-002]
Let \(W\subset V_\Gamma\) be an isotropic stabilizer subspace and let
\(\chi:W\to\{\pm1\}\) be a stabilizer character.  For each finite
\(F\Subset\Gamma\), the coefficient rule
\[
  \omega_F(P_v)=
  \begin{cases}
    \chi(v), & v\in W\cap V_F,\\
    0, & v\notin W\cap V_F
  \end{cases}
\]
defines a positive unital linear functional on \(\cA_F\).  These functionals
are compatible with finite inclusions.
\source{STATE-001 coefficient rule in \path{blueprints/002-canonical-state.tex};
project proof note 004, Sections 2 and 4, especially the product-form
argument after adapted symplectic reduction.}
\lean{Formalization.Stabilizer.FiniteCanonicalState}
\backend{Project-local proof using PAULI-002, SYMP-001 for finite \(F\), and
mathlib positive linear maps if the finite algebra is represented as a
C*-algebra.}
\status{Planned project-local blueprint target.}
\end{theorem}

\begin{proof}[Proof sketch]
Restrict \(W\cap V_F\) to the finite symplectic space \(V_F\) and choose an
adapted finite symplectic basis.  The coefficient rule becomes a product of
one-site pure \(Z\)-eigenstates on fixed modes and normalized traces on
unfixed modes.  Each one-site factor is a positive unital functional; finite
tensor products of states are states.  Compatibility follows because
restricting from \(G\) to \(F\subset G\) drops tensor factors and leaves the
same Pauli coefficients on labels supported in \(F\).
\end{proof}

\begin{theorem}[ID: STATE-003]
The compatible family \((\omega_F)_{F\Subset\Gamma}\) of STATE-002 extends
uniquely to a state \(\omega\) on the quasi-local algebra \(\cA_\Gamma\), or
to the chosen local-net state interface.
\source{Standard C*-inductive-limit state extension principle; for this
project the required compatibility is reduced to the finite-local statement in
STATE-002.  Exact textbook locator still needs page-level extraction before an
Aristotle packet is sent.}
\lean{Formalization.Stabilizer.QuasiLocalCanonicalState}
\backend{Project-local if the quasi-local algebra is abstracted; otherwise use
mathlib positive linear maps and completion/extension lemmas where available.}
\status{Planned project-local blueprint target; source locator needs refinement.}
\end{theorem}

\begin{proof}[Proof sketch]
Define \(\omega\) on the algebraic union by choosing any finite support
containing the observable and applying the corresponding \(\omega_F\).
Compatibility makes this well-defined.  Positivity and norm one hold on every
finite local algebra, hence on the algebraic union.  Continuity extends the
functional uniquely to the norm closure.
\end{proof}

\begin{lemma}[ID: STATE-004]
Let \(S'\) be a consistent stabilizer extension and let \(\Lambda\) be a
region.  The restriction of the canonical state to \(\cA_\Lambda\) has
nonzero Pauli coefficient at a label \(v\in V_\Lambda\) exactly when
\[
  v\in W'\cap V_\Lambda.
\]
Crossing stabilizers affect the restricted state only through products whose
total label lies in \(V_\Lambda\).
\source{Project proof note 004, Lemma 2.4 and Section 2 proof, lines around
the ``Crossing stabilizers'' argument.}
\lean{Formalization.Stabilizer.restrictedCoefficient_eq_internal}
\backend{Project-local theorem over coefficient states.}
\status{Planned project-local blueprint target.}
\end{lemma}

\begin{proof}[Proof sketch]
The global coefficient rule says that \(\omega(P_v)\) is nonzero exactly for
\(v\in W'\).  Since \(P_v\) is supported in \(\Lambda\), the label also lies in
\(V_\Lambda\).  Thus nonzero coefficients are indexed precisely by
\(W'\cap V_\Lambda\).  Stabilizers with support crossing the boundary cannot
contribute unless their product has a label supported entirely in \(\Lambda\).
\end{proof}

\section{Symplectic Normal Form and Product State}

\begin{lemma}[ID: SYMP-001]
Let \(V_\Lambda\) be finite or countable dimensional over \(\Ftwo\), with its
nondegenerate symplectic form, and let \(W\subset V_\Lambda\) be isotropic.
There is a symplectic basis \(\{z_i,x_i:i\in I\}\), with \(I\) finite or
countable, and a subset \(P\subset I\), such that
\[
  W=\operatorname{span}_{\Ftwo}\{z_i:i\in P\}.
\]
\source{Project proof note 004, Lemma 3.1.}
\lean{Formalization.Stabilizer.adaptedSymplecticBasis}
\backend{Project-local algebra over \(\texttt{ZMod 2}\).}
\status{Planned project-local blueprint target.}
\end{lemma}

\begin{proof}[Proof sketch]
Choose a basis of \(W\).  Pair each chosen \(z_i\) with an \(x_i\) by
nondegeneracy, correcting against previously chosen pairs by the usual
finite-stage symplectic Gram--Schmidt procedure.  In the countable case,
iterate over an enumeration.  The span of the chosen hyperbolic pairs has a
symplectic orthogonal complement; choose a countable symplectic basis there
and take the union.
\end{proof}

\begin{lemma}[ID: SYMP-002]
An adapted symplectic basis of \(V_\Lambda\) induces a \(*\)-isomorphism from
the formal qubit quasi-local algebra on the basis index set \(I\) to
\(\cA_\Lambda\), sending formal Pauli labels
\(\sum_{i\in F}(a_i z_i+b_i x_i)\) to the corresponding Pauli operator in
\(\cA_\Lambda\), up to the fixed phase convention.
\source{Project proof note 004, Lemma 3.2.}
\lean{Formalization.Stabilizer.adaptedPauliIso}
\backend{Project-local theorem using finite Pauli algebra universality and
inductive-limit compatibility.}
\status{Planned project-local blueprint target.}
\end{lemma}

\begin{proof}[Proof sketch]
On every finite subset of basis modes the adapted basis is a finite
symplectic basis, so the finite Pauli relations give a unique finite matrix
algebra isomorphism.  These finite isomorphisms are compatible under inclusion
because the phase and generator conventions are fixed.  Passing to the
quasi-local closure gives the required \(*\)-isomorphism.
\end{proof}

\begin{proposition}[ID: PROD-001]
After transporting the restricted canonical state along the isomorphism of
SYMP-002, the state on the formal qubit algebra over \(I\) is
\[
  \bigotimes_{i\in P}\omega_{Z,+}
  \otimes
  \bigotimes_{i\in I\setminus P}\operatorname{tr}_2,
\]
with \(\omega_{Z,+}\) replaced by the corresponding \(Z_i=\pm1\) eigenstate
for a nontrivial stabilizer character.
\source{Project proof note 004, Proposition 4.1.}
\lean{Formalization.Stabilizer.productFormCanonicalState}
\backend{Project-local theorem over Pauli coefficients.}
\status{Planned project-local blueprint target.}
\end{proposition}

\begin{proof}[Proof sketch]
In adapted coordinates, \(v=\sum_{i\in F}(a_i z_i+b_i x_i)\) lies in \(W\)
exactly when all \(b_i=0\) and \(a_i=0\) for \(i\notin P\).  These are exactly
the one-site Pauli coefficients of a pure \(Z\)-eigenstate on \(P\) and the
normalized trace on \(I\setminus P\).  Since Pauli coefficients of finite
tensor products multiply, the local state has the displayed product form.
\end{proof}

\section{Hilbert and GNS Layer}

\begin{definition}[ID: HILB-001]
The Hilbert spaces used in the final proof are:
\begin{enumerate}
\item the GNS Hilbert space \(\cH_\omega\) of a quasi-local state;
\item finite tensor products \(\bigotimes_{i\in F}\mathbb C^2\);
\item the incomplete infinite tensor product
\(\bigotimes_i(\cH_i,\xi_i)\) when importing standard product-factor
identifications;
\item \(\ell^2\) or any fixed separable infinite-dimensional Hilbert space
for type \(I_\infty\) factors.
\end{enumerate}
\source{Mathlib Hilbert-space infrastructure from the coverage audit;
Blackadar III.3.1.1, PDF p. 301, for incomplete infinite tensor products.}
\lean{Formalization.OperatorAlgebra.HilbertInterfaces}
\backend{Mathlib-backed for Hilbert spaces and finite tensor products;
project-local wrappers for incomplete infinite tensor products.}
\status{Planned interface target.}
\end{definition}

\begin{theorem}[ID: GNS-001]
For a state \(\omega\) on a unital C*-algebra \(A\), there are a Hilbert
space \(\cH_\omega\), a unital \(*\)-representation
\[
  \pi_\omega:A\to B(\cH_\omega),
\]
and a cyclic unit vector \(\Omega_\omega\) such that
\[
  \omega(a)=\langle \pi_\omega(a)\Omega_\omega,\Omega_\omega\rangle.
\]
The triple is unique up to unitary equivalence among cyclic triples inducing
\(\omega\).
\source{Kadison--Ringrose I, Theorem 4.5.2 and Proposition 4.5.3, PDF
pp. 287--288; mathlib
\path{Mathlib.Analysis.CStarAlgebra.GelfandNaimarkSegal}.}
\lean{Formalization.GNS.GNSPackage}
\backend{Mathlib-backed through \path{PositiveLinearMap.GNS} and
\path{PositiveLinearMap.gnsStarAlgHom}, with project-local wrappers for
the cyclic vector and reconstruction formula if needed.}
\status{Existing blueprint target; source locator now fixed.}
\end{theorem}

\begin{lemma}[ID: GNS-002]
Let \(\alpha:A\to B\) be a unital \(*\)-isomorphism and let
\(\varphi\) be a state on \(B\).  The GNS representation of
\(\varphi\circ\alpha\) is unitarily equivalent to the pullback of the GNS
representation of \(\varphi\) along \(\alpha\).  Under this unitary, generated
von Neumann algebras are conjugate.
\source{Kadison--Ringrose I, Proposition 4.5.3, PDF pp. 287--288, applied to
the two cyclic triples.}
\lean{Formalization.GNS.gnsUnitaryOfStateIso}
\backend{Project-local wrapper around GNS uniqueness.}
\status{Planned project-local blueprint target.}
\end{lemma}

\begin{proof}[Proof sketch]
The GNS triple of \(\varphi\) composed with \(\alpha\) is a cyclic triple for
\(\varphi\circ\alpha\).  By GNS uniqueness, there is a unitary intertwiner
from the canonical GNS triple of \(\varphi\circ\alpha\) to this pulled-back
triple.  Intertwining on the represented C*-algebra passes to bicommutants by
unitary conjugation.
\end{proof}

\begin{theorem}[ID: GNS-003]
For the product state in PROD-001, the GNS representation is the incomplete
tensor product of the one-site GNS representations, and the represented
von Neumann algebra is the infinite tensor product generated by the one-site
represented algebras with respect to the GNS cyclic vectors.
\source{Blackadar III.3.1.1--III.3.1.3, PDF pp. 301--302.}
\lean{Formalization.GNS.productStateGNS}
\backend{Project-local sourced theorem; mathlib does not currently provide
infinite Hilbert tensor products.}
\status{Planned project-local blueprint target.}
\end{theorem}

\section{Concrete Von Neumann Algebra Layer}

\begin{definition}[ID: VN-001]
Given a state \(\omega\) on \(\cA_\Gamma\), its GNS representation
\(\pi_\omega\), and a region \(\Lambda\subseteq\Gamma\), the represented
region von Neumann algebra is
\[
  M_\Lambda(\omega)=\pi_\omega(\cA_\Lambda)''\subset B(\cH_\omega).
\]
\source{Project proof note 004, Theorem 5.1 statement; mathlib
\path{Mathlib.Analysis.VonNeumannAlgebra.Basic} for concrete
\path{VonNeumannAlgebra H} and commutants.}
\lean{Formalization.OperatorAlgebra.RegionVNPackage}
\backend{Mathlib-backed definition plus project-local packaging from a
subalgebra image to its bicommutant.}
\status{Existing blueprint target.}
\end{definition}

\begin{definition}[ID: VN-002]
For a concrete von Neumann algebra \(M\subset B(\cH)\), define its center by
\[
  Z(M)=M\cap M'
\]
and say that \(M\) is a factor when \(Z(M)=\mathbb C1\).  Equivalently, every
central element of \(M\) is scalar.
\source{Standard von Neumann algebra definition; KR II uses factors in
Corollary 6.5.3, PDF pp. 59--60.  Exact center-definition page locator should
be extracted before packet submission.}
\lean{Formalization.OperatorAlgebra.IsFactor}
\backend{Project-local; mathlib audit found no high-level factor predicate.}
\status{Planned project-local blueprint target; source locator needs refinement.}
\end{definition}

\begin{lemma}[ID: VN-003]
If two represented region algebras are related by unitary conjugation of their
GNS representations, then their bicommutants are also unitarily conjugate and
therefore have the same factor/type classification.
\source{Mathlib concrete commutant infrastructure; GNS uniqueness source KR I,
Proposition 4.5.3, PDF pp. 287--288.}
\lean{Formalization.OperatorAlgebra.unitaryConj_regionVN}
\backend{Mathlib-backed commutant manipulation plus project-local wrapper.}
\status{Planned project-local blueprint target.}
\end{lemma}

\begin{proof}[Proof sketch]
If \(U\pi(a)U^*=\rho(\alpha(a))\) on the represented C*-algebra, then
conjugation by \(U\) sends the commutant of \(\pi(A)\) to the commutant of
\(\rho(\alpha(A))\).  Applying the same identity to commutants gives the
double-commutant statement.  Type predicates are defined in terms of
projections, traces, and centers, all preserved by unitary conjugation.
\end{proof}

\begin{definition}[ID: PROJ-001]
For projections \(p,q\in M\), define Murray--von Neumann equivalence
\(p\sim q\) when there exists a partial isometry \(v\in M\) with
\[
  v^*v=p,\qquad vv^*=q.
\]
A projection \(p\) is finite when no proper subprojection of \(p\) is
equivalent to \(p\), infinite otherwise, and properly infinite when it splits
into two orthogonal subprojections each equivalent to \(p\).
\source{KR II, Chapter 6 comparison theory and Definition 6.5.1, PDF
pp. 57--59, where type definitions are stated in terms of abelian, finite,
and properly infinite projections.  Exact earlier projection-comparison
definition locators still need extraction.}
\lean{Formalization.OperatorAlgebra.ProjectionComparison}
\backend{Project-local; mathlib audit found no Murray--von Neumann API.}
\status{Planned project-local blueprint target; source locator needs refinement.}
\end{definition}

\begin{definition}[ID: TRACE-001]
A trace on a von Neumann algebra \(M\) is a positive linear functional or
weight \(\tau\) satisfying \(\tau(x^*x)=\tau(xx^*)\) on its domain.  The
classification proof needs only the following scoped instances: normalized
finite matrix traces, the tracial product state on the hyperfinite \(II_1\)
factor, and the semifinite tensor trace on \(B(\ell^2)\overline\otimes R\).
\source{Blackadar III.3.1.2--III.3.1.6, PDF pp. 301--304, for product traces
and the \(II_\infty\) model; KR II, Corollary 6.5.3 and Remark 6.5.4, PDF
pp. 59--60, for finite/properly infinite factor distinctions.}
\lean{Formalization.OperatorAlgebra.TraceInterfaces}
\backend{Project-local scoped substitute, not a full normal-weight library
unless later Lean work requires it.}
\status{Planned project-local blueprint target.}
\end{definition}

\begin{definition}[ID: TYPE-DEF-001]
For factors, define the type predicates used by the final theorem as follows.
A factor is type \(I_n\) when its identity is the sum of \(n\) equivalent
minimal projections; type \(I_\infty\) when it is type \(I\) with infinite
identity; type \(II_1\) when it has no minimal projections and the identity is
finite; type \(II_\infty\) when it has no minimal projections and the identity
is infinite but the algebra has nonzero finite projections; and type \(III\)
when it has no nonzero finite projections.
\source{KR II, Definition 6.5.1, Theorem 6.5.2, Corollary 6.5.3, and Remark
6.5.4, PDF pp. 57--60.}
\lean{Formalization.OperatorAlgebra.FactorTypes}
\backend{Project-local; mathlib audit found no type-classification API.}
\status{Planned project-local blueprint target.}
\end{definition}

\begin{theorem}[ID: TYPE-ALT-001]
Every factor has exactly one of the type alternatives needed here:
\[
  I_n,\quad I_\infty,\quad II_1,\quad II_\infty,\quad III.
\]
For the final theorem, it suffices to prove uniqueness and recognition for the
model factors appearing in MATRIX-001, ITP-PURE-001, ITP-TRACE-001, and
TENSOR-001.
\source{KR II, Corollary 6.5.3, PDF pp. 59--60.}
\lean{Formalization.OperatorAlgebra.factor_type_alternatives}
\backend{Project-local theorem, possibly scoped to model factors first.}
\status{Planned project-local blueprint target.}
\end{theorem}

\section{Model Factors and Tensor Type Calculus}

\begin{theorem}[ID: MATRIX-001]
Finite matrix algebras \(M_n(\mathbb C)\) are type \(I_n\) factors with their
normalized traces.  \(B(\ell^2)\), or \(B(\cH)\) for a separable
infinite-dimensional Hilbert space \(\cH\), is a type \(I_\infty\) factor.
\source{KR II, Theorem 6.6.1, PDF pp. 62--63, identifying type \(I_n\)
factors with \(B(\cH)\) for \(\dim\cH=n\); mathlib matrix and trace
infrastructure from the coverage audit.}
\lean{Formalization.OperatorAlgebra.MatrixFactors}
\backend{Project-local theorem using mathlib finite matrices and bounded
operators.}
\status{Planned project-local blueprint target.}
\end{theorem}

\begin{proof}[Proof sketch]
The center of \(M_n(\mathbb C)\) is scalar by elementary matrix-unit
calculus, so it is a factor.  Rank-one matrix units give minimal projections,
and the identity is the sum of \(n\) equivalent minimal projections.  For
\(B(\ell^2)\), rank-one projections are minimal and the identity is a countably
infinite sum of equivalent rank-one projections, giving type \(I_\infty\).
\end{proof}

\begin{theorem}[ID: ITP-001]
Let \((M_i,\varphi_i)\) be finite-dimensional von Neumann algebras with normal
states indexed by a countable set.  Their infinite tensor product in the GNS
representation of the product state is a von Neumann algebra on the incomplete
tensor-product Hilbert space.  If all \(M_i\) are factors, the infinite tensor
product is a factor.  If all \(\varphi_i\) are traces, the product state is a
trace.
\source{Blackadar III.3.1.1--III.3.1.3, PDF pp. 301--302.}
\lean{Formalization.OperatorAlgebra.InfiniteTensorProductFactors}
\backend{Project-local sourced theorem; not expected from mathlib currently.}
\status{Planned project-local blueprint target.}
\end{theorem}

\begin{theorem}[ID: ITP-PURE-001]
The infinite tensor product of one-site pure \(Z\)-eigenstate representations
of \(M_2(\mathbb C)\), over a countably infinite set of fixed modes, generates
\[
  B\!\left(\bigotimes_i^\Omega \mathbb C^2\right),
\]
a type \(I_\infty\) factor.  Over finitely many fixed modes \(p\), the algebra
is \(M_{2^p}(\mathbb C)\), a finite type \(I\) factor.
\source{Blackadar III.3.1.1--III.3.1.3, PDF pp. 301--302, for incomplete
tensor products and factor preservation; KR II, Theorem 6.6.1, PDF pp. 62--63,
for type \(I\) factors as \(B(\cH)\).}
\lean{Formalization.OperatorAlgebra.PureModeProductFactor}
\backend{Project-local sourced theorem.}
\status{Planned project-local blueprint target.}
\end{theorem}

\begin{proof}[Proof sketch]
A pure one-site state gives the defining representation on a one-site Hilbert
space with cyclic vector.  The represented algebra on each finite block is the
full bounded algebra of that finite tensor product.  The infinite incomplete
tensor product is separable and infinite dimensional for countably infinitely
many nontrivial factors.  The generated algebra is therefore \(B(\cH)\), hence
type \(I_\infty\); the finite case is the finite matrix algebra.
\end{proof}

\begin{theorem}[ID: ITP-TRACE-001]
The infinite tracial tensor product
\[
  \overline{\bigotimes}_{i\in\mathbb N}(M_2(\mathbb C),\operatorname{tr}_2)
\]
is the hyperfinite type \(II_1\) factor \(R\).  For finitely many defect modes
\(r\), the tracial product algebra is \(M_{2^r}(\mathbb C)\).
\source{Blackadar III.3.1.4, PDF p. 302, for the tracial \(M_2\) product as
the \(II_1\) factor \(R\); Blackadar III.3.4.1--III.3.4.2, PDF p. 317, for
hyperfinite and infinite tensor products of finite-dimensional algebras.}
\lean{Formalization.OperatorAlgebra.TracialDefectProductFactor}
\backend{Project-local sourced theorem.}
\status{Planned project-local blueprint target.}
\end{theorem}

\begin{proof}[Proof sketch]
The finite case is a finite matrix tensor product with normalized trace.  In
the countably infinite case, Blackadar's infinite tensor product construction
with tracial one-site states gives a factor, and the distinguished product
state is a trace.  The increasing finite matrix tensor subalgebras have
weakly dense union, so the factor is hyperfinite.  The standard example is
the hyperfinite \(II_1\) factor \(R\).
\end{proof}

\begin{theorem}[ID: TENSOR-001]
The tensor products of the pure-mode and defect-mode model factors have the
following types:
\[
\begin{array}{c|c|c}
\hbox{pure factor} & \hbox{defect factor} & \hbox{tensor type}\\
\hline
M_{2^p} & M_{2^r} & I_{2^{p+r}}\\
B(\ell^2) & M_{2^r} & I_\infty\\
M_{2^p} & R & II_1\\
B(\ell^2) & R & II_\infty.
\end{array}
\]
\source{Blackadar III.3.1.6, PDF p. 302, for \(R\overline\otimes B(\cH)\)
as a type \(II_\infty\) factor; KR II, Corollary 6.5.3 and Theorem 6.6.1,
PDF pp. 59--63, for type recognition.}
\lean{Formalization.OperatorAlgebra.TensorProductTypeCalculus}
\backend{Project-local sourced theorem; can initially be scoped exactly to
these four model products.}
\status{Planned project-local blueprint target.}
\end{theorem}

\begin{proof}[Proof sketch]
Tensoring a factor with a finite matrix algebra amplifies multiplicity but
does not change the finite/properly infinite or minimal-projection structure
except for the finite type \(I\) dimension.  \(B(\ell^2)\overline\otimes
M_n\cong B(\ell^2)\), giving type \(I_\infty\).  \(M_n\overline\otimes R\)
is again a finite hyperfinite \(II_1\) factor with normalized trace.  Finally,
\(B(\ell^2)\overline\otimes R\) is the standard \(II_\infty\) amplification
of \(R\).
\end{proof}

\section{Final Classification Targets}

\begin{theorem}[ID: TYPE-001]
Let \(\Lambda\) be finite or countably infinite and let \(S'\) be a consistent
stabilizer extension with canonical projection state \(\omega_{S'}\).  Let
\[
  W'_\Lambda=W'\cap V_\Lambda.
\]
Choose an adapted symplectic basis from SYMP-001, and define
\[
  p=\dim_{\Ftwo}W'_\Lambda,\qquad r=|I\setminus P|.
\]
Then the represented region von Neumann algebra
\[
  M_\Lambda(S')=\pi_{\omega_{S'}}(\cA_\Lambda)''
\]
has type:
\[
\begin{array}{c|c|c}
p<\infty & r<\infty & \hbox{finite type } I\\
p=\infty & r<\infty & \hbox{type } I_\infty\\
p<\infty & r=\infty & \hbox{hyperfinite type } II_1\\
p=\infty & r=\infty & \hbox{type } II_\infty.
\end{array}
\]
\source{Project proof note 004, Theorem 5.1; operator-algebra backend sources
listed in the Source Map above.}
\lean{Formalization.Stabilizer.RegionType}
\backend{Final project theorem after all preceding targets.}
\status{Existing target TYPE-001; master dependency plan complete, proof
blueprint still to be expanded into a dedicated final note.}
\end{theorem}

\begin{proof}[Proof sketch]
By STATE-004, the restricted state sees only the internal stabilizer space
\(W'_\Lambda\).  By SYMP-001 and SYMP-002, after a Pauli algebra isomorphism
the internal stabilizers are exactly the \(Z_i\)-fixed modes for \(i\in P\).
By PROD-001, the transported state is the product of pure one-site states on
\(P\) and tracial one-site states on \(I\setminus P\).  By GNS-002,
VN-003, and GNS-003, the represented region von Neumann algebra is unitarily
equivalent to the tensor product of the pure-mode factor and the defect-mode
factor.  ITP-PURE-001 identifies the pure-mode factor as finite type \(I\) if
\(p<\infty\) and type \(I_\infty\) if \(p=\infty\).  ITP-TRACE-001 identifies
the defect factor as finite matrix if \(r<\infty\) and hyperfinite \(II_1\) if
\(r=\infty\).  TENSOR-001 gives exactly the four displayed cases.
\end{proof}

\begin{corollary}[ID: TYPE-002]
No canonical stabilizer projection state in the countable internal-stabilizer
classification produces a type \(III\) represented region von Neumann algebra.
\source{Project proof note 004, Theorem 5.1 and its final sentence; KR II,
Corollary 6.5.3, PDF pp. 59--60, for the factor type \(III\) criterion.}
\lean{Formalization.Stabilizer.noTypeIII_region}
\backend{Final project corollary.}
\status{Planned project-local blueprint target.}
\end{corollary}

\begin{proof}[Proof sketch]
TYPE-001 identifies every represented region algebra in this sector as one of
finite type \(I\), type \(I_\infty\), hyperfinite \(II_1\), or \(II_\infty\).
By TYPE-DEF-001 and TYPE-ALT-001, type \(III\) factors are precisely the
remaining case with no nonzero finite projections.  Each model factor in the
four rows has nonzero finite projections or a finite/semi-finite trace as
recorded in MATRIX-001, ITP-TRACE-001, and TENSOR-001.  Therefore type \(III\)
does not occur.
\end{proof}

\subsection*{Completion Criterion}

The strengthened blueprint goal is complete when every target in the
dependency table has either:
\begin{enumerate}
\item a dedicated \texttt{blueprints/*.tex} note with source locator,
definition/theorem statement, assumptions, proof sketch or proof, and intended
Lean name; or
\item a section in this master blueprint that is explicitly accepted as the
informal source of truth for that target.
\end{enumerate}
Before any new Lean declaration is introduced, the relevant target must be
moved from ``planned'' to ``blueprint note complete'' in \texttt{MANIFEST.toml}.

\end{document}
